% Generated draft: requires mathematical and source review.
\documentclass[11pt]{article}
\usepackage[margin=25mm]{geometry}
\usepackage{fontspec}
\setmainfont{Latin Modern Roman}
\usepackage{amsmath,amssymb,mathtools,unicode-math}
\setmathfont{Latin Modern Math}
\usepackage{enumitem,needspace,xurl,hyperref}
\setlength{\parindent}{0pt}
\setlength{\parskip}{5pt}
\setlength{\emergencystretch}{3em}
\newcommand{\N}{\mathbb{N}}
\newcommand{\Z}{\mathbb{Z}}
\newcommand{\Q}{\mathbb{Q}}
\newcommand{\R}{\mathbb{R}}
\newcommand{\C}{\mathbb{C}}
\newcommand{\F}{\mathbb{F}}
\newcommand{\A}{\mathbb{A}}
\newcommand{\PP}{\mathbb{P}}
\newcommand{\E}{\mathbb{E}}
\newcommand{\eps}{\varepsilon}
\newcommand{\dd}{\,\mathrm{d}}
\begin{document}
\section*{UnsolvedMath ID: 800214}
\textbf{Source problem code:} catalogue entry
\subsection{Definitions and mathematical statement}
Let \(\mathbb B\subset\C^N\) be the unit ball, where the ambient dimension \(N\) is unspecified, and let \(X\subset\mathbb B\) be a complex analytic subvariety with an isolated singularity at \(0\). Thus \(0\) is singular and some neighborhood of \(0\) in \(X\) contains no other singular points. Let \((\Omega_t)_{t\in I}\) be the source's continuously increasing family of Stein open subsets of \(X\), where \(I\subset\R\) denotes its unspecified parameter range. Increasing means \(\Omega_s\subseteq\Omega_t\) whenever \(s,t\in I\) and \(s<t\). The source does not specify the meaning of continuity for this family; this remains a qualification on the formulation. Assume \(1/2\in I\) and
\[
0\in\partial_X\Omega_{1/2},\qquad
0\in\Omega_t\quad\text{for every }t\in I\text{ with }t>1/2.
\]
Here the boundary is interpreted relative to \(X\), and Steinness and holomorphic functions are understood in the sense of complex spaces, including at singular points.

For open subsets \(U\subseteq V\subseteq X\), say that \(U\) is Runge in \(V\) if for every \(f\in\mathcal O(U)\), every compact \(K\subset U\), and every \(\varepsilon>0\), there exists \(g\in\mathcal O(V)\) such that
\[
\sup_{x\in K}|f(x)-g(x)|<\varepsilon.
\]
The question is whether, under the stated hypotheses,
\[
\Omega_s\text{ is Runge in }\Omega_t
\qquad\text{for every }s,t\in I\text{ with }s<t.
\]
A fully specified statement requires clarification of the parameter range and the continuity condition. The description of the singularity above is local; the source does not expressly state whether other singularities away from the origin are excluded.
\subsection{Short English statement}
Must every earlier member of the given continuously increasing family of Stein open subsets of \(X\) be Runge in every later member, when the singular origin lies on the boundary at parameter \(1/2\) and belongs to every later member? The parameter range and the precise continuity condition remain unspecified.
\subsection{Sources}
\begin{itemize}
\item Local UnsolvedMath registry, permanent ID \(800214\), exact source problem code \texttt{catalogue entry}, title: \emph{Several Complex Variables Problem List --- Problem 10}. The supplied record attributes the problem to Buzzard, Fornaess, and Krantz, with proposed year 1995.
\item Buzzard, Fornaess, and Krantz, \emph{The Several Complex Variables Problem List} (1995), Problem 10. Supplied source URL: \url{https://arxiv.org/abs/math/9509202}. The background records access on 2026-10-16 and extraction from arXiv LaTeX; these are supplied provenance statements, not independently verified access or source checks.
\item The supplied registry labels the problem \emph{open}. This is a supplied status label, not an independently verified assessment of the mathematical literature.
\end{itemize}
\end{document}
